\subsection{A finite-channel theorem for a generic class}
\label{subsec:v45-intro-generic}

The intrinsic inverse has a consequence for a broad class of periodic
tables, without assuming in advance the existence of a suitable measured
graph.  Let $\mathscr P_N^\omega$ be the space of labelled periodic tables
with $N$ connected real-analytic strictly convex obstacle boundaries of
positive curvature and disjoint lifted closures.  The marked rank-two deck
group is fixed; its Euclidean realization is not observed.  Equip support
functions with the relative $C^2$ topology and the lattice matrix with its
usual topology.

\begin{theorem}[Finite-channel periodic rigidity]
\label{thm:v45-intro-finite-channel}
The tables for which every obstacle has trivial proper Euclidean symmetry
form an open dense subset of $\mathscr P_N^\omega$.  For every table in
this subset, there is a selection of $N+1$ marked channels and a positive
offset at which their $N+1$ onsets and $2N+2$ same-type signed limiting
endpoint laws determine the complete labelled table and its Euclidean
period lattice up to one simultaneous proper Euclidean motion.

The selected graph is a spanning tree with two additional edges of
independent deck-cycle gain.  Its labels, gates and rank persist in a
geometric neighborhood.  No common laboratory registration, supplied
curvatures or lattice metric are used in the inverse.
\end{theorem}

The full statement and proof are given in
Theorem~\ref{thm:v45-generic-determination}.  The geometric input is
Theorem~\ref{thm:v45-clear-skeleton}, which applies even to symmetric tables.
If a closest segment between two lifted obstacles meets a third body, that
body splits the pair into two smaller-distance pairs.  Uniform positive
separation makes this descent finite, and its total terminal edge distance
never exceeds the original gap.  Applying it to paths reaching the obstacle
types and two independent deck translates produces a finite clear graph.
A spanning tree and two independent fundamental gains reduce the graph to
$N+1$ channels.  The cycle-space dimension proves that this count is minimal
for the connected two-cycle mechanism, not for arbitrary observation maps.
The selected channels are genuine local alternating billiard channels;
the closed graph walks used for holonomy are not asserted to be specular
multi-channel orbits.

The inverse then follows the previously stated analytic chain: relative
bridge laws, single-offset amplitude cancellation, smooth finite-remainder
contact inversion, and analytic image determination.  Proper obstacle
asymmetry makes every incidence signature unique.  The tree and the two
cycles can therefore be registered in one frame.  This supplies the
architecture required by the more general conditional gluing theorem; it
does not weaken that theorem or replace its local proof by graph algebra.
On an explicit open dense harmonic sublocus, second and third support
coefficients of the \emph{recovered complete images} give all relative
rotations.  Integrating their center-displacement cochain on the tree and
taking its two cycle defects yields both lattice columns and all obstacle
placements; see Proposition~\ref{prop:v45-explicit-image-inverse}.

Genericity refers to the support topology just specified, not to a measure
on tables.  Fixed gates do not reveal exact contacts, and selecting channels
from an unknown unmarked trajectory is a different acquisition problem.
Uniform physical reconstruction is obtained on compact analytic subfamilies
within a persistent design neighborhood, with the same charged pilot and
rare-event caps as before.  Image-registration stability is not promoted to
an effective law-to-entire-boundary rate.  The observations remain different
from the marked length data of \cite{DSKL,FinamoreLeguil}; the theorem is
not a removal of hypotheses from those spectral results.
